% Introduction to the knight vs knight chapter, for the 1500-1599 book only.
% There is no shared texts/knight-vs-knight.tex yet; if one is written, \input it at the end.

\section*{What to look for}
``Knight endgames are pawn endgames'', Botvinnik once said. Dvoretsky explains what he meant: many
laws of pawn endings hold here too, such as the active king, the outside passed pawn, the
breakthrough, shouldering and zugzwang~\cite{Dvoretsky:2008-endgame-manual}. That is true, and you
will meet some of these ideas in this chapter, for example a shoulder check\index{Shoulder check (king)}
in \xref{Diagram}{harizopoulos-klabatseas-54w}.

But the games in this chapter tell a different story. Knight endings are far more spectacular than
pawn endings. They are full of tactics, and long lines in which the defender must find one only move
after another. Most points here were lost to a few tactical ideas.

\emph{Forks and deflections.} The knight forks king and piece, or king and pawn,\index{Knight fork}
and it lures the enemy knight away from a passed pawn by offering itself:\index{Deflection (knight)}
\xref{Diagram}{costa-dasilva-50b} and \xref{Diagram}{dajcman-kerman-66b}. Before every king move,
check which squares the enemy knight can reach. Even mate is possible: see
\xref{Diagram}{vilcacontreras-laurasaavedra-93b}.

\emph{Count the race.} When both sides have passed pawns, count before you act, and ask who queens
with check (\xref{Diagram}{vorniceanu-paraschiv-46b}).\index{Practical play!Counting}

\emph{Unusual ideas save difficult positions.} Promoting to a knight instead of a queen can be the
only way to avoid a fork, or to hold a lost-looking position\index{Underpromotion}
(\xref{Diagram}{evenshaug-syse-56b} and \xref{Diagram}{mitsiou-skodras-56b}). A single knight can
stop two connected pawns\index{Knight against pawns!One knight against connected pawns}
(\xref{Diagram}{krivicius-jeriominas-70b}).

\emph{The defender's simplest resource.} When you are worse, look for a way to exchange or remove
the last pawns\index{Defensive resources!Exchanging or removing the pawns}
(\xref{Diagram}{sabahi-saygili-57b}). When you are better, first stop your opponent's ideas with a
little prophylaxis\index{Strategy!Prophylaxis} (\xref{Diagram}{guedikian-estoyanoff-51w}).
